% Abstract

% ~166 words:
Vessel diameter is a widely used biomarker of vascular health, commonly assessed in the retina using optical coherence tomography angiography (OCTA). However, shear-induced alignment of red blood cells at vessel edges reduces backscattering and suppresses the OCTA signal. This leads to systematic underestimation of vessel diameter. In this study, we established ground-truth diameters in the in vivo mouse retina using Intralipid contrast enhancement, compared binary mask based calculation against model-based profile fitting on maximum and mean intensity projections, and used flicker stimulation to evaluate the performance of each method. Before contrast, a generalised Gaussian fit to the maximum intensity projection retained the strongest linear relationship with ground truth of all methods tested, although it underestimated diameter by approximately 35\%, making it correctable by calibration. Post-contrast, the same fit applied to the mean intensity projection agreed with ground truth ($m = 1.001$, $R^2 = 0.97$). Under flicker stimulation, model-based fitting detected the expected vasodilation, whereas binary mask based calculation instead reported vasoconstriction. These findings characterise OCTA diameter underestimation and establish model-based fitting as a practical means of correcting it.
